\subsection{A persistent obstruction to prescribed coefficient-tail primitives}
\label{sec:B-tail-obstruction}

We record the precise limit of applying the corrected lifting result
to the delayed coefficient modules in SA08, C18.  This gives an actual
admissible target, including a globally free parameter.  It does not
give a counterexample to common sources: the target constructed below
has an explicit admissible common source with the standard theory.

In this subsection $k=\mathbb Z$ and
$S=\Gamma_{\mathbb Z}(T)=\bigoplus_{a\geq0}\mathbb Z T^{[a]}$,
with $\operatorname{wt}(T^{[a]})=a$ and form degree zero.
For $M\geq1$ put
$I_M=\bigoplus_{a\geq M}\mathbb Z T^{[a]}\subset S$.
For a weight-graded $S$-module $V$ write
$(V\langle u\rangle)^w=V^{w-u}$; thus the generator of
$S\langle u\rangle$ has weight $u$.  All maps in the next
proposition are weight-preserving $S$-linear maps.

\begin{proposition}[A-L01: Proved; persistence for the canonical free epimorphism]
\label{AL01}
Let
\[
 F=\bigoplus_{a\geq1} S e_a,\qquad
 \operatorname{wt}(e_a)=a+2,\qquad
 \pi:F\longrightarrow I_1\langle2\rangle,\quad
 e_a\longmapsto T^{[a]} .
 \tag{AL.1}
\]
For every $M\geq1$ there is no map
$s_M:I_M\langle2\rangle\to F$ satisfying
$\pi s_M=\iota_M$, where
$\iota_M:I_M\langle2\rangle\hookrightarrow I_1\langle2\rangle$
is inclusion.  In particular none of these prescribed inclusions
lifts through the displayed free epimorphism, however far the
coefficient tail is delayed.
\end{proposition}

\begin{proof}
The formula for $\pi$ is homogeneous, $S$-linear and surjective.
Suppose that $s_M$ exists.  Reduction modulo a prime $p$ preserves
the alleged map and its equation; the modules $F$ and $I_M$ have the
displayed free abelian bases.  Choose $m=p^r\geq M$, with $r\geq0$.
Homogeneity gives
\[
 \overline{s_M}(T^{[m]})
   =\sum_{a=1}^{m}c_aT^{[m-a]}e_a,\qquad c_a\in\mathbb F_p .
\]
In $\mathbb F_p[X]$ the identity
$(1+X)^{p^r}=1+X^{p^r}$ follows by iterating the characteristic-$p$
binomial identity.  Consequently
$\binom{m}{a}=0$ in $\mathbb F_p$ for $0<a<m$.
Applying $\overline\pi$ to the preceding display therefore gives
$c_mT^{[m]}=T^{[m]}$, so $c_m=1$.

Put $b=(p-1)m$.  The same identity for the power $pm=p^{r+1}$
gives $\binom{pm}{m}=0$ in $\mathbb F_p$.  Hence
$T^{[b]}T^{[m]}=0$ in $I_M/pI_M$.  By $S/pS$-linearity,
$T^{[b]}\overline{s_M}(T^{[m]})$ must be zero.
Its $e_m$ component is $T^{[b]}e_m$, which is nonzero in the free
summand $(S/pS)e_m$.  No summand indexed by a different $a$
contributes to this component.  This is the contradiction.
The argument works with $p=2$ alone and so proves the assertion
over $\mathbb Z$; it also proves its scalar-extended version over
$\mathbb Z_{(p)}$ and $\mathbb Z_p$.
\end{proof}

\begin{theorem}[A-L02: Proved; an admissible target without delayed strict primitives]
\label{AL02}
There is a simplicial, nonnegatively form- and weight-graded,
augmented strict PD differential $S$-algebra $A$ with the following
properties.
\begin{enumerate}
\item Every $A_n^{w,q}$ is a finite free abelian group and
      $A_n^{w,q}=0$ for $q>w$.
\item For $w>q$, the simplicial abelian group $A^{w,q}$ is acyclic,
      including its zeroth homotopy group.
\item For every simplex degree $n$ and every $w\geq3$,
      the coefficient map
      $S^w[0]\to A_n^{w,*}$ is a quasi-isomorphism.
      The bound $3$ works for all $n$ and all form degrees.
\item Write $\Phi_n=\mathbb Z^{n+1}$, with face maps deleting a
      coordinate and degeneracies repeating one.  There is a closed
      simplicial, weight-preserving $S$-linear map
      \[
       \theta:S\langle2\rangle\otimes_{\mathbb Z}\Phi
           \longrightarrow Z^2(A).
       \tag{AL.2}
      \]
      For no $M\geq1$ is its restriction to
      $I_M\langle2\rangle\otimes_{\mathbb Z}\Phi$
      of the form $d_A\beta$, with
      \[
       \beta:I_M\langle2\rangle\otimes_{\mathbb Z}\Phi
                    \longrightarrow A^1
      \]
      a simplicial, weight-preserving $S$-linear map.
\end{enumerate}
Thus eventual vertical exactness and horizontal acyclicity do not
make a prescribed free-parameter cycle admit a compatible
$S$-linear primitive after a coefficient delay.
\end{theorem}

\begin{proof}
Let $K=\ker(\pi)$ for the epimorphism in \textup{(AL.1)}.  Define
the cochain complex of graded $S$-modules
\[
 C^0=K,\qquad C^1=F,\qquad C^2=S\langle2\rangle,
 \qquad d^0=\text{inclusion},\quad d^1=\iota_1\pi ,
 \tag{AL.3}
\]
with all other terms zero.  Its only nonzero cohomology is
$H^2(C)=(S/I_1)\langle2\rangle$, a copy of $\mathbb Z$ in weight
two.  The map $d^0$ is injective, its image is the kernel of $d^1$,
and the image of $d^1$ is $I_1\langle2\rangle$; these statements
prove the cohomology calculation directly.  Each weight piece of
$F$ is finite free; each weight piece of $K$ is a subgroup of a
finite free abelian group and is therefore finite free.

We now specify the standard theory and the comparison maps we use.
The marked constant vector $1=(1,\ldots,1)$ gives a simplicial
monomorphism $\mathbb Z\to\Phi$.  Put
$\Psi=\Phi/\mathbb Z\cdot1$ and
\[
 G=\Gamma_{\mathbb Z}(\Phi)\otimes_{\mathbb Z}\Lambda_{\mathbb Z}(\Psi).
 \tag{AL.4}
\]
Elements of $\Phi$ have weight one and form degree zero; elements
of $\Psi$ have weight one and form degree one.  The differential
sends a generator $\phi\in\Phi$ to its class in $\Psi$ and extends
by the PD derivation rule.  The coefficient map $S\to G$ sends
$T$ to the marked constant vector.  The augmentation kills all
positive weights.  This is the standard elementary theory used
in SA08, C06; the properties needed here follow as follows.

In simplex degree $n$, the marked copy of $\mathbb Z$ is a direct
summand of $\Phi_n$, so a choice of complement identifies
$G_n$, as a weight-preserving PD differential $S$-algebra, with
$S$ tensored with $n$ one-variable divided-power de Rham disks.
A disk has basis $z^{[a]}$ in form degree zero and
$z^{[a]}dz$ in degree one, with
$d(z^{[a+1]})=z^{[a]}dz$.
The map $h(z^{[a]}dz)=z^{[a+1]}$, $h(z^{[a]})=0$ contracts its
positive-weight part.  Tensoring these finite disk contractions
gives an $S$-linear, weight-preserving cochain homotopy equivalence
$G_n\simeq S[0]$.  For precision, the tensor contraction for two
strong deformation retractions $(i,p,h)$ and $(i',p',h')$ is
$h\otimes\mathrm{id}+ip\otimes h'$, with the usual sign
$(f\otimes g)(x\otimes y)=(-1)^{|g||x|}f(x)\otimes g(y)$;
its differential commutator is
$\mathrm{id}\otimes\mathrm{id}-ip\otimes i'p'$.
Induction gives the asserted finite tensor contraction.

The normalized homological complex of $\Phi$ is the identity
disk $\mathbb Z\to\mathbb Z$ in degrees one and zero.  Hence
$\Phi$ has an additive simplicial contraction.
A positive divided-power functor preserves its simplicial
homotopy identities and sends the zero linear map to zero.
Therefore $\Gamma^a(\Phi)$ is simplicially contractible for
$a>0$.  Tensoring a simplicial contraction with a simplicial
module again gives a contraction: on the homotopy maps one uses
the contraction on the first factor and the corresponding
degeneracy on the second.  It follows from
$G^{w,q}=\Gamma^{w-q}(\Phi)\otimes\Lambda^q(\Psi)$ that
$G^{w,q}$ is contractible when $w>q$.  These arguments also
verify that the constructions respect the coefficient map and
the two gradings, without requiring a simplicial choice of the
vertical complements.

Regard $C$ as a constant simplicial cochain complex, and form the
simplicial differential $G$-module
\[
 \mathcal M=G\otimes_S C\otimes_{\mathbb Z}\Phi.
 \tag{AL.5}
\]
The last factor has form and weight degree zero, with zero form
differential.  The form differential is
$d(g\otimes c\otimes\phi)=dg\otimes c\otimes\phi+
(-1)^{|g|}g\otimes dc\otimes\phi$.
The $S$-linear tensor products in \textup{(AL.5)} exist because
the specified coefficient map makes $G$ a differential
$S$-algebra.

Define $A=G\ltimes\mathcal M$, the split square-zero extension:
\[
 (g,m)(g',m')=(gg',gm'+mg'),\qquad
 d_A(g,m)=(d_Gg,d_{\mathcal M}m).
 \tag{AL.6}
\]
Here the right action is the graded-commutative action, so
$mg'=(-1)^{|m||g'|}g'm$.  Associativity and the differential
Leibniz identity follow by expanding both sides and using the
module associativity and differential module identity.
The unit is $(1,0)$.  Every odd square vanishes, because the two
mixed terms have opposite signs and the square of an odd
element of $G$ vanishes.

The positive-weight ideal is the distinguished PD differential
ideal.  For an even homogeneous pair in it put
\[
 \gamma_a(g,m)=
 \bigl(\gamma_a(g),\,\gamma_{a-1}(g)m\bigr)\quad(a\geq1),
 \qquad \gamma_0(g,m)=(1,0).
 \tag{AL.7}
\]
Use the PD sum formula on sums of homogeneous even pairs.
This specifies integral operations of the correct weight and
form degree.  Each module in \textup{(AL.6)} is torsion free:
the explicit simplexwise decomposition of $G_n$ above is free
over $S$, and the terms of $C$ and $\Phi_n$ are free abelian.
Consequently $A\to A\otimes_{\mathbb Z}\mathbb Q$ is injective.
Over $\mathbb Q$, formula \textup{(AL.7)} is exactly
$(g+m)^a/a!$, since $\mathcal M^2=0$.  All PD identities,
the odd-product PD identity, and
$d\gamma_a(z)=\gamma_{a-1}(z)dz$ therefore hold rationally
and hence integrally by this injection.  This proves every
required PD and differential compatibility.  All faces,
degeneracies and the maps $S\to G\to A\to\mathbb Z$ preserve
these operations by their displayed formulas.

For vertical cohomology, tensor the $S$-linear,
weight-preserving contraction $G_n\simeq S$ with $C$ and the
free group $\Phi_n$.  It gives
$\mathcal M_n\simeq C\otimes_{\mathbb Z}\Phi_n$ as cochain
complexes of graded $S$-modules.  Thus its only cohomology is
$H^2(\mathcal M_n)=\Phi_n$ in weight two.
The splitting $A=G\oplus\mathcal M$ is a splitting of
cochain complexes, so $H^0(A_n)=S$ and all additional
cohomology lies in form degree two, weight two.  This proves
assertion (3), with the claimed uniform bound.

For horizontal acyclicity, fix a form and a weight.  The
corresponding component of $\mathcal M$ is a finite direct
sum of simplicial modules of the form
$V\otimes_{\mathbb Z}\Phi$, with $V$ a simplicial module.
The last factor's contraction makes each summand
contractible by the homotopy formula already used.
Hence every row of $\mathcal M$ is contractible.
The calculation for $G$ proves (2).
The finite free bases and finite sums in each fixed weight
prove the finite freeness in (1).
In $C$, the nonzero terms have respectively form degrees
zero, one, two, and weights at least three, three, two.
Thus $q\leq w$ in $C$, and the same inequality in $G$
is preserved by tensor products, proving the support
assertion in (1).

Define \textup{(AL.2)} by
\[
 \theta(c\otimes\phi)=(0,1\otimes c\otimes\phi),
 \qquad c\in C^2=S\langle2\rangle .
 \tag{AL.8}
\]
It is closed because $d_G1=0$ and $d_C(C^2)=0$.
It is simplicial, $S$-linear and weight-preserving by
construction.  Finally, $G_0=S$, concentrated in form
degree zero, and $\Phi_0=\mathbb Z$.  Thus $A_0^1=F$,
$A_0^2=S\langle2\rangle$, and the degree-zero simplicial
component of $\theta$ is the identity of $S\langle2\rangle$.
A map $\beta$ as in (4), evaluated at simplex degree zero,
would therefore be a map
$I_M\langle2\rangle\to F$ lifting its inclusion into
$I_1\langle2\rangle$.  Proposition~\ref{AL01} excludes it.
This proves (4).
\end{proof}

The inclusion $G\to A$, $g\mapsto(g,0)$, and the projection
$A\to G$, $(g,m)\mapsto g$, are actual augmentation- and
coefficient-preserving morphisms of simplicial PD differential
algebras, by \textup{(AL.6)--(AL.7)}; their composite on $G$
is its identity.  Both theories are admissible by the proof
above.  Consequently this example obstructs the specified
primitive-extension method, while explicitly supplying the
common source $G$ for this pair.  It leaves the general
common-source and terminal-localization assertions undecided.

\begin{theorem}[A-L03: Proved; the prescribed first nonlinear PD correction can fail]
\label{AL03}
Over $\mathbb Z$ there is an admissible strict PD theory $A$,
finite free in every simplex, form and weight, with
$A^{w,q}=0$ for $q>w$, and a simplicial element
$x\in A^{2,1}$ such that the following statements hold.
Put $y=dx$ and $\vartheta=\gamma_2(y)\in Z^4(A^{4,*})$.
For every $M\geq1$, the coefficient-linear cycle map
\[
 I_M\langle4\rangle\longrightarrow Z^4(A),\qquad
 T^{[a]}\longmapsto T^{[a]}\vartheta
 \tag{AL.9}
\]
has no simplicial, weight-preserving $S$-linear primitive in $A^3$.
The coefficient module in \textup{(AL.9)} is constant simplicial.
Nevertheless $A_n^{w,*}$ resolves $S^w$ for every $n$ and every
$w\geq5$, with a bound independent of the form degree.
There are actual strict PD morphisms $G\to A\to G$ whose
composite is the identity of the standard theory $G$.

In particular the first new primitive required by the
$p=2$ odd-disk closure in SA08, C18, need not exist with its
prescribed image after any coefficient-tail delay.  The
assertion here concerns that prescribed odd disk; it does
not rule out choosing other maps when constructing a common
source.
\end{theorem}

\begin{proof}
We give all the algebraic operations.  Let
\[
 F' = S u\oplus\bigoplus_{a\geq1}S e_a,\qquad
 \operatorname{wt}(u)=4,\quad
 \operatorname{wt}(e_a)=a+4,
\]
and define the weight-preserving $S$-linear map
\[
 \pi':F'\longrightarrow S\langle4\rangle,\qquad
 \pi'(u)=2,\quad\pi'(e_a)=T^{[a]} .
 \tag{AL.10}
\]
Write $K'=\ker\pi'$.  Give $K'$ form degree two, $F'$
form degree three, and write $S\vartheta$ for
$S\langle4\rangle$ in form degree four.
Let $Sx$ and $Sy$ be free graded $S$-modules with
\[
 ( \operatorname{wt}(x),|x|)=(2,1),\qquad
 ( \operatorname{wt}(y),|y|)=(2,2).
\]
The underlying bigraded $S$-module of our augmented algebra $B$ is
\[
 B=S\oplus Sx\oplus Sy\oplus K'\oplus F'\oplus S\vartheta .
 \tag{AL.11}
\]
The summand $S$ is the coefficient subalgebra.  Multiplication
is $S$-bilinear, its only nonzero products between the five
other summands are
\[
 xy=yx=u,\qquad y^2=2\vartheta ,
 \tag{AL.12}
\]
and $K'\oplus F'\oplus S\vartheta$ annihilates every
positive-form summand.  In particular $x^2=0$.
Every product of three positive-form elements is zero, by
\textup{(AL.12)} and the annihilator rule.  Thus associativity
on three positive-form inputs holds; associativity involving
$S$ is the module associativity.  These exhaust the inputs
by trilinearity.  Koszul commutativity follows from
\textup{(AL.12)} and the zero products, and every odd square
is zero.

Define the $S$-linear differential by
\[
 dx=y,\qquad dy=0,\qquad
 dk=k\in F'\ (k\in K'),\qquad
 df=\pi'(f)\vartheta\ (f\in F'),\qquad d\vartheta=0 .
 \tag{AL.13}
\]
Here $\pi'(f)\vartheta$ denotes the coefficient of
$\pi'(f)\in S\langle4\rangle$ multiplying the indicated
basis symbol.  The equation $d^2k=0$ is precisely the
definition of $K'$, and $d^2=0$ on the other summands is
immediate from the formulas.  The differential preserves
the annihilator summand, so its Leibniz identity on a zero
product involving that summand has both sides zero.
The remaining checks are
\[
 d(x^2)=yx-xy=0,\qquad
 d(xy)=y^2=2\vartheta=du,\qquad
 d(y^2)=0=d(2\vartheta).
 \tag{AL.14}
\]
The check for $yx$ is the same with its displayed even
degree, and scalar inputs use $S$-linearity.  Thus all
products satisfy the differential identity.

The augmentation $B\to S$ is projection, and its kernel
$B_+=Sx\oplus Sy\oplus K'\oplus F'\oplus S\vartheta$
is a differential ideal.  The distinguished PD ideal of
$B$ is its positive-weight ideal, which also contains
$S_+$.  On $S_+$ use its given divided powers.  On the other
homogeneous even elements define
\[
 \gamma_2(a y+k)=a^2\vartheta,\quad
 \gamma_n(a y+k)=0\ (n\geq3),\qquad
 \gamma_n(b\vartheta)=0\ (n\geq2),
 \tag{AL.15}
\]
with $\gamma_1$ the identity and $\gamma_0=1$.
Extend these formulas by the PD sum identity.  More explicitly,
for an even sum $s+z+b\vartheta$, with $s\in S_+$,
$z=ay+k$, and the convention $\gamma_i(s)=0$ for $i<0$, the
formula is
\[
 \gamma_n(s+z+b\vartheta)
 =\gamma_n(s)+\gamma_{n-1}(s)(z+b\vartheta)
       +\gamma_{n-2}(s)a^2\vartheta
 \quad(n\geq1).
 \tag{AL.16}
\]
The displayed formulas are integral and preserve their
specified form and weight degrees.  To check all PD identities,
first observe that every weight piece of $K'$ is a subgroup
of a finite free abelian group, hence free, so
$B\to B\otimes\mathbb Q$ is injective.  Formulas
\textup{(AL.15)--(AL.16)} become exactly $z^n/n!$ in that
rational algebra: $(ay+k)^2=2a^2\vartheta$,
$(ay+k)^3=0$, and $\vartheta^2=0$.
All multiplication, addition, scalar and iteration identities
therefore hold integrally by injectivity.
An odd product is zero in $B$, so its higher divided powers
are zero.  Finally, on $ay+k$ the only nontrivial higher
differential check is
$d\gamma_2(ay+k)=0=(ay+k)dk$, by the annihilator rule.
On $S_+$ and $S\vartheta$ it is immediate, and the sum formula
and Leibniz rule give it on every even sum.  This verifies
the strict PD differential structure in full.

As a cochain complex, $B$ is the direct sum of $S[0]$, the
identity disk $Sx\to Sy$, and
$K'\hookrightarrow F'\xrightarrow{\pi'}S\vartheta$.
The latter complex has cohomology only in form degree four,
where it is
\[
 (S/(2,I_1))\langle4\rangle
       =(\mathbb Z/2)\langle4\rangle .
 \tag{AL.17}
\]
Indeed the image of $\pi'$ is exactly $(2,I_1)$, and its
kernel is $K'$.  This proves the cohomology assertion
without changing the coefficient algebra.

We next make the target simplicial while preserving the
displayed nonlinearity.  Use the standard $G$ and $\Phi$
from the proof of Theorem~\ref{AL02}.  In this paragraph
$\Phi_n=\mathbb Z^{n+1}$ also has its coordinatewise
multiplication and marked unit $(1,\ldots,1)$.
Deleting or repeating coordinates is a ring map, so this
additional multiplication is compatible with exactly the
same faces and degeneracies used for its additive contraction.
Put $R_n=G_n\otimes_S B$.  The augmentation of $B$ gives
a split PD differential augmentation $R_n\to G_n$.
Here is a direct verification of the PD structure on the tensor
product.  It is torsion free because $G_n$ is free over $S$.
In the rational tensor product use $z^a/a!$ for even $z$.
The integral tensor product is closed under these operations:
for an even pure tensor with both factors even and its
$B$-factor in $B_+$, the formula is
$\gamma_a(g\otimes b)=g^a\otimes\gamma_a(b)$;
for a pure tensor of two odd factors its square and all
higher divided powers are zero.  The base summand is $G_n$
with its existing operations.  The PD sum formula treats
every finite sum of such tensors.  These expressions agree
rationally and therefore integrally whenever two presentations
of the same element are used.  The rational identities and
injectivity prove all PD identities and differential
compatibility, as before.  This supplies the asserted
strict PD tensor product and its augmentation.

Define the simplicial PD differential algebra by the fiber
product
\[
 A_n=\underbrace{R_n\times_{G_n}\cdots\times_{G_n}R_n}_{n+1
             \text{ factors}}
     =G_n\oplus (G_n\otimes_S B_+)^{n+1}.
 \tag{AL.18}
\]
All operations in the first description are coordinatewise.
The second description identifies its ideal with
$G\otimes_S B_+\otimes_{\mathbb Z}\Phi$.
Faces apply the corresponding face of $G$ and delete the
same numbered coordinate; degeneracies apply the
corresponding degeneracy of $G$ and repeat that coordinate.
They preserve the common image in $G_n$, so these formulas
define the fiber product simplicially.  They commute with
products, divided powers and differential, since the
operations are coordinatewise.  The coefficient map is
the diagonal copy of $S$ in $G$, and the augmentation to
$\mathbb Z$ is the augmentation of $G$.  Its
positive-weight ideal is differential and PD stable.
In particular no new weight-zero coefficients or form
generators have been introduced.

Every weight-form component of the added ideal is a
simplicial module $V\otimes\Phi$ and is therefore
simplicially contractible.  The calculation for $G$ gives
horizontal acyclicity of $A^{w,q}$ whenever $w>q$.
Vertically, the weight-preserving $S$-linear contraction
$G_n\simeq S$ gives
$G_n\otimes_SB_+\simeq B_+$.
Equation \textup{(AL.17)} consequently gives
$H^0(A_n)=S$, with the sole additional cohomology
$(\mathbb Z/2)^{n+1}$ in form degree four and weight four.
This proves the uniform vertical resolution assertion.
All weight pieces of $B$ are finite free, so the corresponding
pieces of $A_n$ are finite free as well.  The generators
$x,y,u,\vartheta$ have the listed weights; all $e_a$ have
weight at least five, and $K'$ has no weight-four part
because multiplication by $2$ on $\mathbb Z u$ is injective.
Thus $q\leq w$ on $B$ and hence on $A$.

In \textup{(AL.18)} take $x$ to be the tuple with each entry
$1\otimes x$ and common image zero in $G_n$.
These tuples are compatible with faces and degeneracies.
Equations \textup{(AL.12)--(AL.16)}, applied coordinatewise,
give $y=dx$, $\vartheta=\gamma_2(y)$ and $du=2\vartheta$.
Thus there is an actual morphism from the full odd disk
$S\otimes\Lambda(x)\otimes\Gamma(y)$, with $dx=y$, to $A$;
its divided powers of $y$ of index at least three map to zero.

Finally $G_0=S$ and there is only one factor in
\textup{(AL.18)}, so $A_0=B$.
A primitive of \textup{(AL.9)}, evaluated in simplex degree
zero, would give a homogeneous $S$-linear map
$I_M\langle4\rangle\to F'$ whose composite with $\pi'$
is the inclusion into $S\langle4\rangle$.
Reduce modulo two and choose $m=2^r\geq M$.
The coefficient of $e_m$ in the image of $T^{[m]}$ must
be one: the $u$ contribution projects to zero modulo two,
and all the intermediate binomial coefficients vanish,
as in Proposition~\ref{AL01}.
Multiplication by $T^{[m]}$ kills $T^{[m]}$ in the source
modulo two, since $\binom{2m}{m}=0$ modulo two, but leaves
the nonzero component $T^{[m]}e_m$ in $F'/2F'$.
This contradiction proves the claimed failure for every $M$.
The diagonal section $G\to A$ and the common augmentation
$A\to G$ in \textup{(AL.18)} are the promised strict PD maps
and their composite is the identity.
\end{proof}

